\section{The structural theorem}
We use the twelve-page preprint arXiv:2407.10388v1 by Hou, Feng,
Leemans and Qu~\cite{HFLQ}. We recover its Theorems 1.1--1.3
and the explicit construction in Theorem 5.1, with the correction stated
below. The results remain attributed to those authors. This edition
reorganizes the accepted proofs through the common paper~\cite{UF};
no previous combined manuscript is an additional proof premise.

We use the conventions of that paper: marked generators are distinct
nonidentity involutions, the endpoints commute, and the adjacent
orders are at least three. A string C-group satisfies the full subset
intersection property. Its type is $\{k_1,k_2\}$, and it is tight
when $|G|=2k_1k_2$. Conjugation is $x^g=g^{-1}xg$.

% BEGIN F1.Structure
\begin{theorem}[The odd kernel]\label{F1.Structure}
Let $p$ be an odd prime and $m\geq2$. A nondegenerate rank-three
string C-group $(G;r_0,r_1,r_2)$ of order $4p^m$ has a normal
Sylow $p$-subgroup $P$, and
\[
 G=P\rtimes K,\qquad K=\gen{r_0,r_2}\cong C_2\times C_2.
\]
Both type parameters are divisible by $p$, at least one is divisible
by $2p$, and $P$ needs exactly two generators.
\end{theorem}
\begin{proof}
The number $n_p$ of Sylow $p$-subgroups divides four and is congruent
to one modulo $p$. The only case other than $n_p=1$ is
$p=3$, $n_p=4$. Assume this case occurs. Conjugation on the four
Sylow subgroups gives a homomorphism $G\longrightarrow S_4$ with
kernel $N\leq N_G(P)=P$, since the normalizer has index four.
The image has order $4\cdot3^j$ dividing $24$, so $j=0$ or $1$.
If $j=0$, then $N=P$ and $P$ is normal, a contradiction. If $j=1$,
the image is the index-two subgroup $A_4$ of $S_4$. Its involutions
all lie in the normal Klein four subgroup. They cannot generate
$A_4$, whereas an image of an involution-generated group is generated
by involutions and identities. This is also a contradiction.
Thus $P$ is normal.

The endpoint subgroup $K$ has order four, since its two generators
are distinct commuting involutions. It intersects the odd group $P$
trivially, and $|PK|=|G|$, giving the stated semidirect product.
Each adjacent dihedral group has order $2k_i$, so
$k_i\mid2p^m$. Nondegeneracy implies $p\mid k_i$.
If both $k_i$ were odd, their images in the exponent-two quotient
$G/P=K$ would force $r_0P=r_1P=r_2P$, contrary to the distinct
images of $r_0,r_2$. Thus one label is even.

The inverted-orbit lemma of~\cite{UF} gives two generators for $P$.
If $P$ were cyclic, the nontrivial $p$-parts of both adjacent rotation
groups would contain its unique subgroup of order $p$. That subgroup
would lie in $G_{\{0,1\}}\cap G_{\{1,2\}}=\gen{r_1}$,
which has order two. This contradiction proves that two generators
are necessary.
\end{proof}
% END F1.Structure
